% !TEX program = lualatex
\documentclass[aspectratio=43,professionalfonts,handout]{beamer}
\usepackage{beamer_template}

\newcommand{\LectureNum}{10}
\newcommand{\LectureTitle}{一般の周期関数のフーリエ級数（2）}
\newcommand{\TextPages}{p.\,84-86}
\newcommand{\CourseName}{応用数学}
\renewcommand{\TermName}{前期}

\begin{document}

\begin{frame}{\CourseName\quad 第\LectureNum 回「\LectureTitle」}
  {\small \TermName\quad \TeacherName\quad ／\quad 教科書 \TextPages}

  \textbf{目標}：フーリエ級数の収束定理を一般の周期関数に適用し、数学的公式を導くことができる

\end{frame}

\begin{frame}{収束定理の復習と応用}
  \begin{block}{}
  \begin{itemize}
    \item 区分的に滑らかな条件
    \item 不連続点での収束値
    \item 具体的な点を代入 $ \to $ 数学的公式
  \end{itemize}
  \end{block}
\end{frame}

\begin{frame}{公式の導出}
  \begin{block}{}
  \begin{itemize}
    \item $x=1/2$ 代入 $ \to 1-1/3+1/5-... = \pi/4$
    \item $x=0$ 代入 $ \to 1/1^{2}+1/3^{2}+1/5^{2}+... = \pi^{2}/8$
    \item 美しい数学的関係
  \end{itemize}
  \end{block}
\end{frame}

\begin{frame}{演習（練習問題1）}
  \begin{exampleblock}{自分で解いてみよう}
  \begin{itemize}
    \item p.90 問1〜4より選題
  \end{itemize}
  \end{exampleblock}
\end{frame}

\begin{frame}{まとめ}
  \begin{itemize}
    \item 収束定理の応用
    \item 次回：複素フーリエ級数
  \end{itemize}
\end{frame}

\end{document}
